\section{Catastrophic Forgetting Mitigation}

Training on new domains risks degrading performance on earlier ones.  The
pipeline employs four complementary strategies.

\subsection{Replay Buffers}

Each stage loads pre-cached chunks from earlier corpora as a replay pool.  At
each batch construction, with probability \texttt{replay\_frac} a training
chunk is replaced with one sampled uniformly from the replay pool.

\textbf{Configuration:}
\begin{itemize}[nosep]
  \item Initial replay fraction: 15\% of samples.
  \item Replay cap: 35--40\% (stage-dependent).
  \item EMA-based adaptive scaling: if TinyStories forgetting exceeds
        thresholds, \texttt{replay\_frac} is increased by
        $\text{scale} \times \text{forgetting\_ema}$ (scale $= 3.0$).
  \item Replay sources accumulate: Stage~3 replays from TinyStories and
        ROCStories; Stage~6 replays from TinyStories and Children-Stories.
\end{itemize}

\subsection{Anchor Validation Baseline}

Before training begins at each stage, the TinyStories validation loss is
computed and stored as an \emph{anchor value}.  Throughout training,
forgetting is defined as:
\[
\text{forgetting} = \frac{L_\text{ts}^\text{current} - L_\text{ts}^\text{anchor}}
                         {\max(L_\text{ts}^\text{anchor},\; 10^{-6})}
\]
This relative metric is tracked both raw and as an exponential moving average
(EMA, $\alpha = 0.1$), and logged at every evaluation step.

\subsection{Synaptic Intelligence (SI)}

SI maintains per-parameter importance estimates accumulated during training.
A quadratic penalty is added to the loss:
\[
\mathcal{L}_\text{SI} = \lambda_\text{SI} \sum_k \Omega_k (\theta_k - \theta_k^*)^2
\]
where $\Omega_k$ is the importance of parameter $k$ and $\theta_k^*$ is its
value at the end of the previous stage.  $\lambda_\text{SI}$ ranges from
0.015 to 0.02 across stages.  The $\Omega$ matrices are saved alongside
checkpoints (\texttt{si\_omega\_stage\_N.pt}).

\subsection{Progressive Layer Thaw}

During architectural expansions (Stages~4 and 6), the bottom 3 layers are
initially frozen.  If forgetting exceeds a threshold (5\%), layers are
progressively unfrozen from bottom to top.  This ensures that early
representations---which encode fundamental linguistic features---are
preserved during the initial adaptation phase.

\subsection{Measuring Forgetting}

Forgetting is monitored through:
\begin{enumerate}[nosep]
  \item \texttt{ts\_forgetting}: raw relative forgetting vs.\ anchor.
  \item \texttt{ts\_forgetting\_ema}: smoothed EMA of forgetting.
  \item Cross-domain val loss: all 7 val sets are evaluated at every
        checkpoint, regardless of the active training domain.
\end{enumerate}
\vspace{2pt}
The final Stage~6 model achieves TinyStories val loss of 1.669
(anchor: 1.734), corresponding to \textbf{negative forgetting} ($-3.8\%$),
indicating the model actually \emph{improved} on the source domain.
